\documentclass[12pt]{article}

\usepackage[margin=1in]{geometry}
\usepackage{enumitem}
\usepackage{amsmath}
\usepackage{array}
\usepackage{titlesec}

\title{Lecture 3 Notes: Aristotle's \textit{On Interpretation}}
\author{}
\date{}

\begin{document}

\maketitle

\section*{Central Theme}

Aristotle's \textit{On Interpretation} explains how spoken and written signs
express thought, and how thoughts become true-or-false propositions through
affirmation and denial.

\[
\boxed{
\textit{On Interpretation} = \text{how words become propositions}
}
\]

In the \textit{Categories}, Aristotle classified the kinds of things that can be
said. In \textit{On Interpretation}, he explains how saying becomes judging.

\[
\text{Categories} \rightarrow \text{what can be said}
\]

\[
\textit{On Interpretation} \rightarrow \text{how sayings become true or false}
\]

\section*{1. The Place of \textit{On Interpretation} in Aristotle's Logic}

\textit{On Interpretation} belongs to Aristotle's logical writings. It follows
naturally after the \textit{Categories}.

The \textit{Categories} dealt with simple terms such as substance, quantity,
quality, relation, place, time, action, and affection. \textit{On Interpretation}
now asks how terms are combined into statements.

\subsection*{Teaching Point}

Aristotle is moving from classification to judgment.

\[
\text{term} \rightarrow \text{sentence} \rightarrow \text{proposition}
\]

Only propositions can be true or false.

\subsection*{Whiteboard}

\[
\boxed{
\text{Logic begins when meaningful words become truth-bearing statements.}
}
\]

\section*{2. Words, Thoughts, and Things}

Aristotle opens with a famous account of how language signifies reality.

\[
\text{written words} \rightarrow \text{spoken words}
\rightarrow \text{mental experiences} \rightarrow \text{things}
\]

Written words symbolize spoken words. Spoken words symbolize mental
experiences. Mental experiences are images of things.

\subsection*{Important Distinction}

Different peoples may use different writing systems and different speech
sounds, but Aristotle says that the mental experiences signified by speech are
the same for all, as are the things of which these experiences are images.

\subsection*{Teaching Point}

Language is conventional, but it is not arbitrary in the deepest sense. Speech
points to thought, and thought points to reality.

\subsection*{Whiteboard}

\[
\boxed{
\text{Language signifies thought; thought represents things.}
}
\]

\[
\text{writing} \rightarrow \text{speech} \rightarrow \text{thought} \rightarrow \text{being}
\]

\section*{3. Isolated Terms Are Not True or False}

Aristotle insists that isolated thoughts or isolated expressions are not yet
true or false.

Examples:

\[
\text{man}
\]

\[
\text{white}
\]

These words signify something, but they do not yet affirm or deny anything.

Even a strange expression such as:

\[
\text{goat-stag}
\]

has significance, but it is not true or false unless something is added, such
as ``is'' or ``is not.''

\subsection*{Teaching Point}

Truth and falsity require combination and separation.

\[
\text{``man''} \quad \text{not true or false}
\]

\[
\text{``man is white''} \quad \text{true or false}
\]

\[
\text{``man is not white''} \quad \text{true or false}
\]

\subsection*{Whiteboard}

\[
\boxed{
\text{Truth and falsity arise only when something is affirmed or denied.}
}
\]

\section*{4. Nouns}

Aristotle next defines the noun.

\subsection*{Definition}

A noun is a significant sound by convention, without reference to time, and no
part of it is significant separately from the whole.

\[
\boxed{
\text{noun} = \text{significant by convention, without time}
}
\]

\subsection*{Examples}

\[
\text{man}
\]

\[
\text{horse}
\]

\[
\text{health}
\]

\subsection*{By Convention}

Aristotle says that nothing is a noun by nature. A sound becomes a noun when it
functions as a symbol by convention.

This means that words signify because a linguistic community uses them that
way.

\subsection*{Simple and Composite Nouns}

Some nouns are simple. Others are composite.

For example:

\[
\text{pirate-boat}
\]

is composite. Its parts contribute to the meaning of the whole, but the part
does not necessarily function independently in the same way.

\subsection*{Indefinite Nouns}

Aristotle treats expressions such as:

\[
\text{not-man}
\]

as indefinite nouns. They are not ordinary nouns, but they still denote
something in an indefinite way.

\section*{5. Verbs}

A verb is significant, like a noun, but it also carries time.

\subsection*{Definition}

\[
\boxed{
\text{verb} = \text{significant by convention, with time}
}
\]

A verb is also a sign of something said of something else.

\subsection*{Example}

\[
\text{health} = \text{noun}
\]

\[
\text{is healthy} = \text{verb}
\]

The difference is that ``is healthy'' indicates the present existence of the
state in question.

\subsection*{Teaching Point}

A verb introduces time and predication.

It allows us to say that something belongs to something, or that something does
not belong to something.

\subsection*{Whiteboard}

\[
\text{noun: health}
\]

\[
\text{verb: is healthy}
\]

\[
\boxed{
\text{The verb makes predication temporal.}
}
\]

\section*{6. Sentences}

A sentence is a significant portion of speech.

But not every sentence is a proposition.

\subsection*{Sentence}

\[
\boxed{
\text{sentence} = \text{meaningful speech}
}
\]

A sentence has meaning, but it may not be true or false.

\subsection*{Example}

\[
\text{May Socrates recover!}
\]

This is meaningful, but it is not true or false. It is a prayer or wish, not a
proposition.

\subsection*{Teaching Point}

Aristotle is not studying every kind of meaningful speech. He is interested in
the kind of speech that can bear truth and falsity.

\section*{7. Propositions}

A proposition is a sentence that is either true or false.

\[
\boxed{
\text{proposition} = \text{meaningful speech that is true or false}
}
\]

Examples:

\[
\text{Socrates is healthy.}
\]

\[
\text{Socrates is not healthy.}
\]

These are propositions because they affirm or deny something.

\subsection*{Simple and Composite Propositions}

A simple proposition asserts or denies one thing of one thing.

A composite proposition is made up of simpler propositions.

\subsection*{Whiteboard}

\[
\text{sentence} \supset \text{proposition}
\]

\[
\text{All propositions are sentences, but not all sentences are propositions.}
\]

\section*{8. Affirmation and Denial}

Aristotle defines the two basic forms of proposition.

\subsection*{Affirmation}

An affirmation is a positive assertion of something about something.

\[
\text{Socrates is white.}
\]

\subsection*{Denial}

A denial is a negative assertion of something about something.

\[
\text{Socrates is not white.}
\]

\subsection*{Teaching Point}

Affirmation and denial are the fundamental operations of truth-bearing speech.

\[
\boxed{
\text{A proposition either affirms or denies.}
}
\]

\section*{9. Contradiction}

Every affirmation has an opposite denial, and every denial has an opposite
affirmation.

A pair of propositions is contradictory when the same predicate is affirmed and
denied of the same subject.

\subsection*{Example}

\[
\text{Socrates is white.}
\]

\[
\text{Socrates is not white.}
\]

These are contradictories.

\subsection*{Important Condition}

The subject and predicate must be the same, and they must not be used
equivocally.

\[
\boxed{
\text{Contradiction requires the same subject and the same predicate.}
}
\]

\subsection*{Teaching Point}

Contradiction is not merely disagreement. It is a precise logical opposition.

\section*{10. Universal and Individual Subjects}

Aristotle distinguishes universal and individual subjects.

\subsection*{Universal}

A universal is something that can be predicated of many subjects.

\[
\text{man}
\]

\[
\text{animal}
\]

\subsection*{Individual}

An individual is something that is not predicated of many subjects.

\[
\text{Callias}
\]

\[
\text{Socrates}
\]

\subsection*{Whiteboard}

\[
\text{universal} = \text{predicable of many}
\]

\[
\text{individual} = \text{not predicable of many}
\]

\section*{11. Universal, Particular, Contrary, and Contradictory Propositions}

Aristotle distinguishes several forms of proposition.

\subsection*{Universal Affirmation}

\[
\text{Every man is white.}
\]

\subsection*{Universal Denial}

\[
\text{No man is white.}
\]

\subsection*{Particular Affirmation}

\[
\text{Some men are white.}
\]

\subsection*{Particular Denial}

\[
\text{Not every man is white.}
\]

or:

\[
\text{Some men are not white.}
\]

\subsection*{Contraries}

Contrary propositions cannot both be true, but they may both be false.

\[
\text{Every man is white.}
\]

\[
\text{No man is white.}
\]

These cannot both be true. But they can both be false, since some men may be
white and some may not be white.

\[
\boxed{
\text{Contraries cannot both be true, but may both be false.}
}
\]

\subsection*{Contradictories}

Contradictory propositions cannot both be true and cannot both be false.

\[
\text{Every man is white.}
\]

\[
\text{Not every man is white.}
\]

One must be true, and the other must be false.

\[
\boxed{
\text{Contradictories cannot both be true and cannot both be false.}
}
\]

\subsection*{Classroom Example}

Original proposition:

\[
\text{Every student is prepared.}
\]

Contradictory:

\[
\text{Not every student is prepared.}
\]

Contrary:

\[
\text{No student is prepared.}
\]

\subsection*{Teaching Point}

The sentence ``not every student is prepared'' does not mean the same thing as
``no student is prepared.''

\[
\boxed{
\text{Not every} \neq \text{none}
}
\]

\section*{12. Single and Compound Propositions}

Aristotle is careful about what makes a proposition genuinely one.

An affirmation or denial is single if it indicates one fact about one subject.

\subsection*{Example of a Single Proposition}

\[
\text{Socrates is white.}
\]

\subsection*{Problem Case}

Some propositions look single but are really compound.

For example, if a word had two unrelated meanings, a proposition using that word
might not express one unified assertion.

\subsection*{Teaching Point}

A proposition is not made one merely by being written as one sentence. It is one
only if it signifies one fact.

\[
\boxed{
\text{Grammatical unity is not always logical unity.}
}
\]

\section*{13. Future Contingents}

Chapter 9 contains one of the most famous arguments in Aristotle's logical
works.

The issue is whether statements about the future are already determinately true
or false.

\subsection*{Present and Past Propositions}

Statements about what is or what has already happened must be either true or
false.

\[
\text{Socrates is sitting.}
\]

\[
\text{Socrates sat yesterday.}
\]

\subsection*{Future Propositions}

Statements about future contingent events are different.

Example:

\[
\text{There will be a sea-fight tomorrow.}
\]

\[
\text{There will not be a sea-fight tomorrow.}
\]

Aristotle says that a sea-fight must either happen tomorrow or not happen
tomorrow, but it is not necessary that it happen, and it is not necessary that
it fail to happen.

\subsection*{Whiteboard}

\[
\boxed{
\text{Either it will happen or it will not happen.}
}
\]

But not:

\[
\boxed{
\text{It is already necessary which one will happen.}
}
\]

\subsection*{Why This Matters}

If every future-tense proposition were already determinately true or false in
the same way as present and past propositions, then everything would seem to
happen by necessity.

That would threaten:

\begin{itemize}[itemsep=0.25em]
    \item deliberation,
    \item action,
    \item chance,
    \item contingency,
    \item responsibility.
\end{itemize}

\subsection*{Teaching Point}

Aristotle preserves contingency by distinguishing between logical alternatives
and necessary outcomes.

\[
\boxed{
\text{The future is open where there is real potentiality in contrary directions.}
}
\]

\section*{14. Logical Necessity and Future Necessity}

Aristotle does not deny that one of the alternatives must occur.

\[
\text{Either there will be a sea-fight or there will not be a sea-fight.}
\]

This is necessary as a disjunction.

But Aristotle denies that either side is already necessary by itself.

\[
\text{It is not necessary that there will be a sea-fight.}
\]

\[
\text{It is not necessary that there will not be a sea-fight.}
\]

\subsection*{Teaching Point}

There is a difference between saying:

\[
\text{Necessarily, either } P \text{ or not-}P
\]

and saying:

\[
\text{Either necessarily } P \text{ or necessarily not-}P
\]

\[
\boxed{
\text{Necessity of the alternative is not necessity of either alternative.}
}
\]

\section*{15. Modal Propositions}

The later chapters examine propositions involving possibility, impossibility,
contingency, and necessity.

\[
\text{possible}
\]

\[
\text{impossible}
\]

\[
\text{contingent}
\]

\[
\text{necessary}
\]

\subsection*{Basic Modal Distinctions}

\[
\text{It may be} \neq \text{It must be}
\]

\[
\text{It may not be} \neq \text{It cannot be}
\]

\subsection*{The Placement of Negation}

Aristotle is especially concerned with finding the correct contradictories of
modal propositions.

For example, the contradictory of:

\[
\text{It may be}
\]

is not:

\[
\text{It may not be}
\]

because the same thing may both be possible and possible not to be.

The proper contradictory is:

\[
\text{It cannot be}
\]

\subsection*{Whiteboard}

\[
\text{It may be} \quad \text{vs.} \quad \text{It cannot be}
\]

\[
\text{It is necessary} \quad \text{vs.} \quad \text{It is not necessary}
\]

\subsection*{Teaching Point}

Modal logic requires special care about where the negation belongs.

\[
\boxed{
\text{Changing the place of ``not'' changes the logic of the proposition.}
}
\]

\section*{16. Possibility and Necessity}

Aristotle distinguishes kinds of possibility.

Something may be possible because it is already actual.

\[
\text{A walking man can walk because he is walking.}
\]

Something else may be possible because it could be actual under the right
conditions.

\[
\text{A seated man can walk because he has the capacity to walk.}
\]

\subsection*{Teaching Point}

Possibility can mean actuality or capacity. Aristotle is careful not to treat
every use of ``possible'' as identical.

\[
\boxed{
\text{Possible is said in more than one way.}
}
\]

\section*{17. The Contrary of an Affirmation}

The final chapter asks whether the contrary of an affirmation is another
affirmation or a denial.

For example, what is the contrary of:

\[
\text{Every man is just}
\]

Is it:

\[
\text{No man is just}
\]

or:

\[
\text{Every man is unjust}
\]

\subsection*{Teaching Point}

The final chapter returns to the question of logical opposition.

Aristotle wants to distinguish carefully between:

\begin{itemize}[itemsep=0.25em]
    \item contradiction,
    \item contrariety,
    \item affirmation,
    \item denial,
    \item contrary predicates.
\end{itemize}

\subsection*{Whiteboard}

\[
\text{Every man is just}
\]

\[
\text{No man is just}
\]

\[
\text{Every man is unjust}
\]

\[
\boxed{
\text{Logical opposition has structure.}
}
\]

\section*{18. Aristotle's Method in \textit{On Interpretation}}

Throughout the treatise, Aristotle follows a steady method.

\[
\boxed{
\text{define} \rightarrow \text{distinguish} \rightarrow \text{test oppositions}
}
\]

He defines basic terms:

\[
\text{noun}, \quad \text{verb}, \quad \text{sentence}, \quad \text{proposition}
\]

He distinguishes kinds of propositions:

\[
\text{affirmation}, \quad \text{denial}, \quad \text{contrary}, \quad \text{contradictory}
\]

He tests difficult cases:

\[
\text{future contingents}, \quad \text{modal propositions}, \quad \text{contrary affirmations}
\]

\section*{19. Conclusion}

\textit{On Interpretation} is Aristotle's bridge from language to logic.

It explains how meaningful words become truth-bearing propositions, and how
propositions stand in relations of affirmation, denial, contradiction,
contrariety, possibility, and necessity.

\[
\boxed{
\textit{On Interpretation} = \text{Aristotle's theory of truth-bearing speech}
}
\]

\section*{Final Framing}

The lecture should leave students with four main ideas:

\begin{enumerate}[itemsep=0.5em]
    \item Words signify thoughts, and thoughts represent things.
    \item Isolated terms are meaningful, but not true or false.
    \item Truth and falsity arise through affirmation and denial.
    \item Contradiction, contrariety, possibility, and necessity must be
    carefully distinguished.
\end{enumerate}

\[
\boxed{
\text{In the \textit{Categories}, Aristotle classified what can be said.}
}
\]

\[
\boxed{
\text{In \textit{On Interpretation}, Aristotle explains how what is said becomes true or false.}
}
\]

\end{document}